%=====================================================================================
\section{Reliability across prompts and models}\label{sec:reliability}

\subsection{Measures}

I measure agreement across the five prompts in two ways. The intraclass correlation ICC(3,1) of
\citet{shrout1979} treats the prompts as fixed raters and asks what share of the variance in a single
prompt's scores is variance between statements, after removing each prompt's average level. It equals 1
when all prompts rank and space the statements identically, up to a constant shift. Krippendorff's alpha
with the interval metric \citep{hayes2007,krippendorff2018} also counts a constant shift between prompts
as disagreement, so it is at or below the ICC when prompts differ in level.

Each measure is computed four ways: on raw score levels; on levels after the rate decision is regressed
out (``resid''); on hold meetings only; and on meeting-to-meeting changes. For ``resid'' each prompt's
scores are regressed on a single code for the decision ($+1$ hike, $0$ hold, $-1$ cut). Separate hike
and cut dummies give similar results (below).
The \emph{drop} is $1-\text{resid}/\text{raw}$. It asks how much of the agreement is agreement about the
decision itself, which no one needs an LLM to read.

\subsection{Results}

Table~\ref{tab:reliability} and Figure~\ref{fig:icc} report the results. The first block crosses the same
\R{icc.fable.n} policy statements from 2015--2019 with all five prompts and four models. On raw levels
the three inexpensive models reach an ICC between \R{icc.haiku.raw} and \R{icc.gem.raw}. Once the rate
decision is removed, their ICC falls to between \R{icc.haiku.resid} and \R{icc.gem.resid}, a drop of
\R{icc.mini.drop} to \R{icc.haiku.drop}. Fable starts at \R{icc.fable.raw} and keeps \R{icc.fable.resid}
(a drop of \R{icc.fable.drop}). The block contains \R{block.hikes} hikes and \R{block.cuts} cuts. With
separate hike and cut dummies instead of the single code, the drops are \R{dummy.cheap.min} to
\R{dummy.cheap.max} for the inexpensive models and \R{dummy.fable.drop} for Fable. Alpha tells the same
story. Its drops for the inexpensive models are
smaller (\R{alpha.haiku.drop} to \R{alpha.gem.drop}) because alpha already penalises their level
differences in the raw scores, but on hold meetings it shows how little some of them agree: alpha is
\R{alpha.haiku.hold} for Claude 3 Haiku and \R{alpha.mini.hold} for GPT-4o mini, against
\R{alpha.fable.hold} for Fable. On meeting-to-meeting changes, the quantity that market regressions use,
the inexpensive models reach an ICC of \R{icc.haiku.chg} to \R{icc.gem.chg} and Fable \R{icc.fable.chg}.

\begin{table}[t]
\centering
\caption{Agreement across five prompts}\label{tab:reliability}
\footnotesize
\setlength{\tabcolsep}{4pt}
% generated by paper/make_tables.py - do not edit
\begin{tabular}{@{}llr ccccc ccccc@{}}
\toprule
 & & & \multicolumn{5}{c}{ICC(3,1)} & \multicolumn{5}{c}{Krippendorff's $\alpha$} \\
\cmidrule(lr){4-8}\cmidrule(l){9-13}
Model & Window & $n$ & Raw & Resid & Hold & Chg & Drop & Raw & Resid & Hold & Chg & Drop \\
\midrule
\multicolumn{13}{@{}l}{\emph{Crossed block: same statements, five prompts, four models}}\\
Gemini 2.5 Flash & 2015-01--2019-12 & 40 & 0.810 & 0.562 & 0.580 & 0.697 & 31\% & 0.791 & 0.564 & 0.524 & 0.698 & 29\% \\
Claude 3 Haiku & 2015-01--2019-12 & 40 & 0.609 & 0.333 & 0.420 & 0.419 & 45\% & 0.398 & 0.336 & 0.084 & 0.422 & 16\% \\
GPT-4o mini & 2015-01--2019-12 & 40 & 0.728 & 0.517 & 0.537 & 0.577 & 29\% & 0.634 & 0.519 & 0.355 & 0.579 & 18\% \\
Claude Fable 5.1 & 2015-01--2019-12 & 40 & 0.957 & 0.929 & 0.937 & 0.920 & 3\% & 0.953 & 0.930 & 0.928 & 0.921 & 2\% \\
\addlinespace
Gemini 2.5 Flash & 2004-01--2007-12 & 33 & 0.912 & 0.830 & 0.907 & 0.696 & 9\% & 0.891 & 0.831 & 0.885 & 0.697 & 7\% \\
\addlinespace\multicolumn{13}{@{}l}{\emph{Full corpus (statements with all five prompts)}}\\
Gemini 2.5 Flash & 2000-02--2026-07 & 221 & 0.934 & 0.874 & 0.885 & 0.729 & 6\% & 0.933 & 0.874 & 0.881 & 0.729 & 6\% \\
Claude Fable 5.1 & 2000-02--2026-07 & 174 & 0.976 & 0.957 & 0.961 & 0.896 & 2\% & 0.972 & 0.957 & 0.950 & 0.896 & 2\% \\
\bottomrule
\end{tabular}

\tabnote{\emph{Notes:} Policy corpus; statements with all five prompts. Raw = score levels; Resid =
levels after regressing each prompt's scores on the rate decision coded $+1/0/-1$; Hold = hold meetings
only; Chg = changes between consecutive corpus statements; Drop $=1-\text{Resid}/\text{Raw}$. ICC(3,1)
follows \citet{shrout1979}; alpha uses the interval metric.}
\end{table}

\begin{figure}[t]
\centering
% generated by paper/make_tables.py
\begin{tikzpicture}
\begin{axis}[ybar, bar width=13pt, width=0.95\linewidth, height=6.2cm,
  ymin=0, ymax=1.12, ytick={0,0.25,0.5,0.75,1}, ylabel={ICC(3,1) across five prompts}, symbolic x coords={Gemini 2.5 Flash,Claude 3 Haiku,GPT-4o mini,Claude Fable 5.1},
  xtick=data, enlarge x limits=0.15, legend style={at={(0.5,-0.14)}, anchor=north, legend columns=-1, font=\footnotesize, draw=none, /tikz/every even column/.append style={column sep=8pt}},
  legend cell align=left, nodes near coords, nodes near coords style={font=\scriptsize, /pgf/number format/fixed, /pgf/number format/precision=2}, ymajorgrids, major grid style={black!15}]
\addplot[fill=black!80, draw=black] coordinates {(Gemini 2.5 Flash,0.810) (Claude 3 Haiku,0.609) (GPT-4o mini,0.728) (Claude Fable 5.1,0.957)};
\addlegendentry{Levels (raw)}
\addplot[fill=black!35, draw=black] coordinates {(Gemini 2.5 Flash,0.562) (Claude 3 Haiku,0.333) (GPT-4o mini,0.517) (Claude Fable 5.1,0.929)};
\addlegendentry{Levels after removing the rate decision}
\addplot[pattern=north east lines, draw=black] coordinates {(Gemini 2.5 Flash,0.697) (Claude 3 Haiku,0.419) (GPT-4o mini,0.577) (Claude Fable 5.1,0.920)};
\addlegendentry{Meeting-to-meeting changes}
\end{axis}
\end{tikzpicture}

\caption{Agreement across prompts by model, 2015--2019 block ($n=\R{icc.fable.n}$ statements). Part of
the inexpensive models' agreement on levels is agreement on the rate decision.}\label{fig:icc}
\end{figure}

The size of the drop depends on the sample. For Gemini it is \R{icc.gem.drop} in the 2015--2019 block but
\R{icc.gem04.drop} on 2004--2007 and \R{icc.gemfull.drop} on the full corpus. ICC is a ratio of
between-statement to total variance, so it is high when the statements in the sample differ a lot in
stance, as they do across the whole period from 2000 to 2026. Within a few years of similar statements,
prompt noise weighs more. The reliability of an LLM score is therefore a property of the model, the
prompt set and the sample together, and a researcher should report it for the sample actually used.

Fable's consistency is also high in a second sense: the \R{retest.n} cells it scored twice in separate
runs correlate at \R{retest.r}.

\subsection{Model choice versus prompt choice}

Prompt sensitivity is only one source of arbitrariness. A researcher also chooses a model. On the crossed
block I decompose the variance of the scores into statement, prompt and model effects and their
interactions, with a random-effects analysis of variance. With one score per cell, the three-way
interaction cannot be separated from noise, so it is part of the residual. Table~\ref{tab:variance}
reports the shares.

\begin{table}[t]
\centering
\caption{Variance decomposition: statement $\times$ prompt $\times$ model, 2015--2019 block}\label{tab:variance}
\small
% generated by paper/make_tables.py - do not edit
\begin{tabular}{@{}llr rrrrrrr@{}}
\toprule
Data & Scaling & $n$ & Statement & Prompt & Model & S$\times$P & S$\times$M & P$\times$M & Residual \\
\midrule
Levels & raw & 40 & 61.0 & 4.5 & 0.0 & 3.0 & 9.7 & 5.7 & 16.1 \\
Levels & std.\ within model & 40 & 61.3 & 4.6 & 0.0 & 3.0 & 9.2 & 6.0 & 15.9 \\
Changes & raw & 39 & 49.1 & 0.0 & 0.0 & 6.7 & 15.2 & 0.0 & 29.0 \\
Changes & std.\ within model & 39 & 50.2 & 0.0 & 0.0 & 6.4 & 15.2 & 0.0 & 28.3 \\
\bottomrule
\end{tabular}

\tabnote{\emph{Notes:} Shares of total variance (\%), random-effects ANOVA, \R{var.n} statements $\times$ 5
prompts $\times$ 4 models (Gemini 2.5 Flash, Claude 3 Haiku, GPT-4o mini, Claude Fable 5.1); negative
variance estimates are set to zero. ``Std.\ within model'' z-scores each model's scores first, which
removes scale differences between models. S$\times$P: prompts disagree about which statements are more
hawkish; S$\times$M: models disagree; the residual contains the three-way interaction and noise.}
\end{table}

On levels, agreed differences between statements account for \R{var.levels.raw.stmt} of the variance.
On changes, the agreed part falls to \R{var.changes.raw.stmt}. The main disagreement on changes is between
models: the statement-by-model interaction, which measures models disagreeing about which statements
became more hawkish, takes \R{var.changes.raw.sxm}, against \R{var.changes.raw.sxp} for the
statement-by-prompt interaction. The residual rises from \R{var.levels.raw.resid} on levels to
\R{var.changes.raw.resid} on changes. Standardising each model's scores first leaves these shares almost
unchanged, so they do not reflect differences in scale. For meeting-to-meeting changes, the choice of
model matters about twice as much as the choice of prompt. A robustness check across prompts alone
understates how much an LLM stance measure depends on the researcher's choices.

\subsection{Is the frontier model consistent because it recognises the text?}

One reading of Fable's high reliability is that it is simply a better reader. Another is that it
recognises each statement and recalls a fixed view of it, whatever the prompt
(Section~\ref{sec:recognition} shows it does recognise them). The two readings predict different
behaviour on statements the model has not seen. Fable's dating probe (Section~\ref{sec:cutoff}) suggests
that its knowledge fades from March 2026, so I scored the seven latest statements with all five prompts
and split them by whether the probe dated them exactly and with high confidence. The prompt spread (the standard deviation of the five scores) averages \R{spread.unknown}
on the \R{spread.unknown.n} statements it probably has not seen (March to July 2026) and \R{spread.known}
on the \R{spread.known.n} it dates confidently, against a median of \R{spread.pre.median} and a 90th
percentile of \R{spread.pre.p90} on the \R{spread.pre.n} pre-2024 statements with all prompts (Appendix
Table~\ref{tab:spread}). All seven recent statements lie between the \R{spread.pctmin} and
\R{spread.pctmax} quantiles of the pre-2024 spread, so Fable is somewhat less consistent on recent
statements than on old ones, whether it knows them or not. Two of the four statements it probably has
not seen are the unverified texts of June and July 2026 (Section~\ref{sec:data}); without them, the mean
spread of the remaining \R{spread.unknown.verified.n} is \R{spread.unknown.verified}, slightly above that
of the known statements. Statements this few cannot settle the question.
